\cotefile{157-1}
\scanpage{212}
\mspage{66}

Rappelons que les \deleted{flèches} \inserted{classes $\mathrm W_{\mathrm{asph}}$} \unsure{$W$-asphériques|homosphériques} et celles \inserted{$\mathrm{Liss}_{W}$} des flèches $W$-lisses, se déterminent mutuellement, ainsi

\[
\left\{\begin{aligned}
(f:X\to Y)\unsure{\in|\notin}\mathrm W_{\mathrm{asph}}&\Longleftrightarrow
  \forall\,(Y'\xrightarrow{u}Y)\in\mathrm{Liss}_{W},\\
  &\quad X'=X\times_Y Y'\longrightarrow X'\ \text{est}\ \in W\\
(\unsure{u|f}:Y'\to Y)\in\mathrm{Liss}_{W}&\Longleftrightarrow
  \text{le foncteur }X\mapsto X'=X\times_Y Y'\\
 &\quad\text{de }\mathrm{Cat}_{/Y}\text{ dans }\mathrm{Cat}_{/Y'}\ \text{transforme}\\
 &\quad\text{flèches dans }\mathrm W_{\mathrm{asph}}\text{ dans}\\
 &\quad\text{flèches dans }\mathrm W_{\mathrm{asph}}
\end{aligned}\right.
\tag*{(4.10)}
\]

Pour toute partie $A$ de $\mathrm{Fl}(\mathrm{Cat})$, \inserted{stable par isom.,} on peut associer une partie $\mathbb{L}_{\deleted{\unsure{W}}}(A)$ de $\mathrm{Fl}(\mathrm{Cat})$, par la relation

\[
\begin{aligned}
(u:Y'\to Y)\in\mathbb{L}_{\deleted{W}}(A)&\Longleftrightarrow u^*:\mathrm{Cat}_{/Y}\longrightarrow\mathrm{Cat}_{/Y'}\\
&\quad\text{transforme flèches}\\
&\quad\text{dans }A\text{ en flèches }\in A.
\end{aligned}
\tag*{(4.11.)}
\]

On voit que $\mathbb{L}_{\deleted{W}}(A)$ contient les iso, est stable par composition, \deleted{mais il n'est pas clair} \deleted{qu'il soit} \inserted{et est} stable aussi par chgt. de base (comme c'est le cas de $\mathrm{Liss}_{W}$). \deleted{\unsure{A d'|à d}} Inversement, à une partie $L$ de $\mathrm{Fl}(\mathrm{Cat})$, stable par isom., on associe une partie $\mathbb A_{\deleted{\unsure{W}}}(L)$ par

\[
\begin{aligned}
(f:X\to Y)\in\mathbb A_{\deleted{W}}(L)&\Longleftrightarrow\forall\,(Y'\to Y)\text{ dans }L,\\
&\quad X'\unsure{\Rightarrow|=}X\times_Y Y'\to Y'\\
&\quad\text{est dans }\deleted{\unsure{A|L}}W
\end{aligned}
\tag*{(4.12)}
\]

Si $L$ est stable par chgt. de base, alors $\mathbb A_{\deleted{W}}(L)$ est stable par composition. Il contient les iso \deleted{\unsure{si $L$ les}} (car $W$ les contient), \deleted{Mais il} \inserted{\unsure{est|et} contenu dans $W$,} \inserted{si $L$ contient les identités}. Mais
% archive sheet number: 212
